\documentclass[10pt,a4paper]{amsart}
\usepackage[margin=19mm]{geometry}
\usepackage{amsmath,amssymb,mathtools}
\usepackage[scr=rsfs,cal=euler]{mathalfa}
\usepackage{xfrac}
\usepackage[unicode,hidelinks,bookmarks=false]{hyperref}
\newcommand{\ie}{i.e.\ }
\newcommand{\cf}{cf.\ }
\newcommand{\resp}{respectively\ }
\newcommand{\Subsection}{\subsection}
\providecommand{\nobreakdash}{\nobreak\hbox{-}}
\setlength{\emergencystretch}{2em}
\multlinegap0pt
\usepackage{bm}
\usepackage{bbm}
\usepackage{dsfont}
\usepackage{xspace}
\usepackage{cancel}
\usepackage{mathtools}
\usepackage[shortlabels]{enumitem}
\usepackage{amsthm}
\makeatletter
\@ifundefined{thm}{\newtheorem{thm}{Thm}}{}
\@ifundefined{claim}{\newtheorem{claim}[thm]{Claim}}{}
\@ifundefined{lem}{\newtheorem{lem}[thm]{Lemma}}{}
\makeatother
\makeatletter
\newcommand{\F}{\mathbb{F}}
\expandafter\ifx\csname poc\endcsname\relax
\newenvironment{poc}{\begin{proof}[Proof of the claim]\renewcommand*{\qedsymbol}{$\blacksquare$}}{\end{proof}}
\fi
\makeatother
\title[Mutual position of two smooth quadrics over finite fields]{Mutual position of two smooth quadrics over finite fields}
\author{Shamil Asgarli, Chi Hoi Yip}
\date{Source: arXiv:2404.06754v2 (CC BY 4.0)}
\begin{document}
\maketitle

\noindent\emph{Source and licence.} Mutual position of two smooth quadrics over finite fields. Version arXiv:2404.06754v2 (CC BY 4.0), distributed under the Creative Commons Attribution 4.0 International licence, as recorded in the arXiv licence metadata for this exact version and shown on the abstract page https://arxiv.org/abs/2404.06754v2. Full rights notice: licenses/arxiv-2404.06754-CC-BY-4.0.txt.\par\smallskip

\noindent\textit{Editorial scope.} Counted unit: the Lem whose source label is \texttt{lem1} in the article (source0.tex lines 319--324, source label \texttt{lem1}), delivered together with the article's own complete original proof of it (source0.tex lines 326--370, one \texttt{proof} environment of 45 lines). No mathematics is written, repaired or re-derived: statement and proof are the article's own text line for line. Required context retained uncounted: lem:RL (lines 302--307). Every definition, display and label of those lines is retained unchanged. Omitted from this package and not redistributed: 1--301, 308--318, 325--325, 371--440. Nothing inside a retained excerpt is omitted or altered: every retained body is delivered exactly as the article's lines are written, so no omission receipt of any kind accompanies this package. Citations: the retained excerpts carry no citation command at all, so no bibliography entry of the article is transcribed and no reference is redistributed. Reference adaptation: none was needed and none was made; every reference inside the delivered unit resolves inside it, so no unresolved reference or citation is produced. Printed slips kept literal: no printed word, symbol, hypothesis or number of the counted statement or of its proof was altered. Layout and class: the article's own class is \texttt{amsart} with packages including fullpage, times, colonequals, amsmath, amssymb, amsthm, url, inputenc, babel, comment, bbm, enumerate, bm, graphicx; the wrapper is the lane's pinned \texttt{amsart} editorial wrapper at 10pt on A4, which restates the theorem-like environments the retained text needs and adds the article's own macro definitions that the retained text needs (\texttt{\textbackslash F}), with extra wrapper packages (bm, bbm, dsfont, xspace, cancel, mathtools, enumitem). The delivered numbering is the unit's own. Assets: no retained span contains a figure environment, a graphics-inclusion command, a PDF-page inclusion, a table or a TikZ picture; no image file is copied into this package, the package declares no asset dependency and no image file is redistributed with it. Licence: version arXiv:2404.06754v2 of this article is distributed under the Creative Commons Attribution 4.0 International licence, as recorded in the arXiv licence metadata for this exact version and shown on the abstract page https://arxiv.org/abs/2404.06754v2. A full rights notice is delivered at licenses/arxiv-2404.06754-CC-BY-4.0.txt. Characters: every retained excerpt is pure ASCII, so the delivered engine, XeLaTeX with its default Latin Modern text font, reports no missing character.
\begin{lem}[Rojas-Le\'{o}n]\label{lem:RL}
 Let $X$ be a projective Cohen-Macaulay scheme defined over $\F_q$ with pure dimension $m \geq 2$, embedded in $\mathbb{P}^N$ as the closed subscheme defined by $r$ homogeneous forms $F_1, \ldots, F_r$ of degrees $a_1, \ldots, a_r$. Let $H$ and $Z$ be homogeneous forms in $\F_q\left[X_0, \ldots, X_N\right]$ of degrees $d$ and $e$, where $\gcd(d, e)=1$, $\gcd(e, p)=1$ with $p$ the characteristic of $\F_q$. We will also denote by $H$ and $Z$ the hypersurfaces they define in $\mathbb{P}^N$. Assume that $X \cap H \cap Z$ has pure codimension 2 in $X$, and denote by $\delta$ the dimension of the singular locus of $X \cap H \cap Z$. Let $\chi$ be a nontrivial multiplicative character of $\F_q$, such that $\chi^e$ is nontrivial. Let $V=X-(H \cup Z)$ and $f\colon V \rightarrow \F_q^*$ be the map defined by $f(\mathbf{x})=H(\mathbf{x})^e / Z(\mathbf{x})^d$. The following estimate holds:
$$
\bigg|\sum_{\mathbf{x} \in V(\F_q)} \chi(f(\mathbf{x}))\bigg| \leq 3(3+\max \left(a_1, \ldots, a_r, e\right)+d)^{N+r+2} \cdot q^{(m+\delta+2) / 2} .
$$   
\end{lem}

\begin{lem}\label{lem1}
Let $q \geq 7$ be an odd prime power and $n \geq 3$ odd. Assume that $C:f=0$ and $D:g=0$ are two distinct smooth quadrics in $\mathbb{P}^{n-1}$, defined over $\F_q$. Then
$$
\left|\sum_{P\in\mathbb{P}^{n-1}(\F_q)} \chi(f(P) g(P))\right| \leq 3\cdot 8^{n+1} q^{(2n-3)/2} +2q^{n-2}. 
$$ 
\end{lem}

\begin{proof} We apply Lemma~\ref{lem:RL} by choosing $X$, $H$, and $Z$ as follows:
\begin{enumerate}[(i)]
\item $X=\mathbb{P}^{n-1}$, $m=n-1$, $N=n-1$, and $r=0$;
\item $H(\mathbf{x})=f(\mathbf{x})g(\mathbf{x})$, and $d=4$;
\item $Z$ is a hyperplane $(e=1)$ which is neither tangent to $C=\{f=0\}$, neither tangent to $D=\{g=0\}$, and nor $Z$ contains any irreducible component of $C\cap D$. 
\end{enumerate}   
We justify why such $Z$ exists. The number of tangent hyperplanes defined over $\mathbb{F}_q$ to either $C$ or $D$ is at most $2(q^{n-2}+q^{n-3}+\cdots+1)$. Moreover, $C\cap D$ has at most $4$ distinct irreducible components, each with dimension $n-3$. The number of hyperplanes in $\mathbb{P}^{n-1}$ containing a linear component (of dimension $n-3$) is $q+1$; if the component is not linear, there is at most one hyperplane in $\mathbb{P}^{n-1}$ containing it. Thus, the number of ``bad" hyperplanes is at most $2(q^{n-2}+\cdots + 1) + 4(q+1)$. However, the total number of hyperplanes is $q^{n-1}+q^{n-2}+\cdots+1$, which is strictly bigger than $2(q^{n-2}+\cdots + 1) + 4(q+1)$ when $q\geq 7$ and $n \geq 3$. In fact, when $n\geq 5$, the intersection $C\cap D$ does not contain a hyperplane component because $C$ is parabolic (the maximum projective dimension of a linear space within $C$ is $\frac{n-3}{2}<n-3$); in this case, the number of ``bad" hyperplanes can be reduced from $2(q^{n-2}+\cdots + 1) + 4(q+1)$ to $2(q^{n-2}+\cdots + 1) + 2$. As a result, one can check that the hyperplane $Z$ exists when $q\geq 3$ and $n\geq 5$.

Observe that $H(\mathbf{x})=f(\mathbf{x})g(\mathbf{x})$ defines a quartic hypersurface. For one of the hypotheses, $X\cap H\cap Z = H\cap Z$ has pure codimension $2$ in $X=\mathbb{P}^{n-1}$ since a hyperplane section of a quartic is a quartic of one less dimension (hence pure codimension $2$ in $X$). Since the hyperplane $Z\cong\mathbb{P}^{n-2}$ does not contain any components of $C\cap D$, the dimension of $C\cap D\cap Z$ is $\delta = n-4$. Indeed, after a projective linear change of variables, the hyperplane $Z$ is defined by the equation $\{x_{n-1}=0\}$; then, $C_0 = C\cap Z$ and $D_0 = D\cap Z$ are two quadrics of one less dimension (with the same equations as $C$ and $D$ after substituting $x_{n-1}=0$). Since $Z$ is tangent to neither $C$ nor $D$, we have that $C_0$ and $D_0$ are smooth quadrics. Moreover, as $Z$ does not contain any irreducible component of $C\cap D$, it follows that $C_0\neq D_0$, and so $C_0\cap D_0 = C\cap D\cap Z$ has the expected dimension of $n-4$.

\begin{claim}
The singular locus of $X\cap H\cap Z$ is equal to $C\cap D\cap Z$; in particular, the dimension of the singular locus of $X\cap H\cap Z$ is $\delta = n-4$.
\end{claim}
\begin{poc}
By abuse of notation, we write $Z(x_0, \ldots, x_{n-1})=0$ for the defining equation of the hyperplane $Z$. The singular locus of $X\cap H\cap Z = H\cap Z$ is the set of points $P$ in $H\cap Z$ such that the following $2\times n$ Jacobian matrix has rank less than $2$:
$$
\operatorname{Jac}_{P}(f(\mathbf{x})g(\mathbf{x}), Z(\mathbf{x})) = 
\begin{pmatrix}
 f(P) \frac{\partial g}{\partial x_0}(P) + g(P) \frac{\partial f}{\partial x_0}(P) &  \cdots &  f(P) \frac{\partial g}{\partial x_{n-1}}(P) + g(P) \frac{\partial f}{\partial x_{n-1}}(P) \\ 
 \frac{\partial Z}{\partial x_0}(P) & \cdots & \frac{\partial Z}{\partial x_{n-1}}(P)
\end{pmatrix}.
$$
Since $P\in H\cap Z$, it follows that $f(P)=0$ or $g(P)=0$. We consider the following three cases. 
\begin{enumerate}
    \item If $f(P)=0$ and $g(P)=0$, then the first row vanishes, and the matrix has rank $1$. 
\item If $f(P)=0$ and $g(P)\neq 0$, then the first row is a nonzero scalar multiple of the gradient vector $\nabla f (P)=\left(\frac{\partial f}{\partial x_0}(P), \ldots, \frac{\partial f}{\partial x_{n-1}}(P) \right)$. In this case, the first row is not a multiple of the second row due to the condition (iii), as the hyperplane $Z$ is not tangent to the quadric $C=\{f=0\}$. Hence, the matrix has full rank (i.e., rank $2$) in this case.
\item If $f(P)\neq 0$ and $g(P)=0$, the same conclusion as in (2) holds by symmetry. 
\end{enumerate}
Therefore, the singular locus of $H\cap Z$ is precisely $C\cap D\cap Z$.
\end{poc}

By Lemma~\ref{lem:RL}, we have
\begin{equation}\label{eq:-Z}
\left|\sum_{P\in\mathbb{P}^{n-1}(\F_q)\setminus Z} \chi(f(P) g(P))\right|=\left|\sum_{P\in\mathbb{P}^{n-1}(\F_q)\setminus (H \cup Z)} \chi\bigg(\frac{f(P) g(P)}{Z(P)^4}\bigg)\right| \leq 3\cdot 8^{n+1} q^{(2n-3)/2}.  
\end{equation}
On the other hand, we have the trivial upper bound 
\begin{equation}\label{eq:Z}
\left|\sum_{P\in Z(\F_q)} \chi(f(P) g(P))\right| \leq q^{n-2}+q^{n-3} + \cdots + q + 1 < 2 q^{n-2}.
\end{equation}
Combining inequality~\eqref{eq:-Z} and inequality~\eqref{eq:Z}, we get
$$
\left|\sum_{P\in\mathbb{P}^{n-1}(\F_q)} \chi(f(P) g(P))\right| \leq 3\cdot 8^{n+1} q^{(2n-3)/2} +2q^{n-2}, 
$$
as required.
\end{proof}

\end{document}
